\documentclass[UTF8]{ctexart}

\usepackage[a4paper,margin=1in]{geometry}
\usepackage{amsmath,amssymb,mathtools,bm}
\usepackage{booktabs,longtable,array,tabularx,makecell}
\usepackage{hyperref}
\usepackage{xurl}
\usepackage{enumitem}
\usepackage{float}

\hypersetup{
    colorlinks=true,
    linkcolor=black,
    urlcolor=blue,
    citecolor=black
}

\setlist[itemize]{leftmargin=2em}
\setlist[enumerate]{leftmargin=2em}
\setlength{\LTleft}{0pt}
\setlength{\LTright}{0pt}
\setlength{\emergencystretch}{3em}
\renewcommand{\arraystretch}{1.16}
\allowdisplaybreaks
\sloppy

\newcolumntype{Y}{>{\raggedright\arraybackslash}X}
\newcommand{\code}[1]{\texttt{\detokenize{#1}}}
\newcommand{\codecell}[1]{\parbox[t]{\linewidth}{\raggedright\ttfamily\small\detokenize{#1}}}
\newcommand{\wrapcell}[1]{\parbox[t]{\linewidth}{\raggedright #1}}
\newcommand{\sg}{\operatorname{sg}}
\newcommand{\rms}{\operatorname{RMS}}

\title{\code{mra.py} 中异常分数与 \code{shift score} 模块的科研式说明手册}
\author{}
\date{2026-04-17}

\begin{document}

\maketitle

\begin{abstract}
本文面向仓库当前版本的 \code{mra.py}，以科研论文说明书的形式分析异常分数计算模块，重点展开其中的 \code{shift score}。文档首先定位代码入口和训练入口，然后说明训练流程、参数更新过程、训练集估计出的参数如何作用于测试集。核心结论是：\code{shift score} 不是一个可反向传播的神经网络层，而是一个基于训练集补全窗口统计特征的分布偏移度量；它先从训练集补全窗口中估计特征均值和标准差，再用这些训练集统计量去衡量测试补全窗口离训练分布有多远。为了让数学原理易懂，本文把公式和代码逐行对应，并给出输入、输出张量形状和数据在内部流动的过程。图学习、频域插补和 Transformer 检测器只在解释 \code{shift score} 上下游依赖时简要说明，不作为本文主体。
\end{abstract}

\noindent\textbf{关键词：} 异常检测，分布偏移，shift score，多采样率，补全窗口，标准化距离

\tableofcontents
\newpage

\section{文档范围与核心结论}

\subsection{本文重点}

本文重点分析以下代码区域：
\begin{itemize}
    \item \code{mra.py:740--773}：各类异常分数的训练集归一化与最终合成；
    \item \code{mra.py:815--843}：\code{shift score} 的特征构造与分布偏移分数计算；
    \item \code{mra.py:993--1055}：训练集和测试集统计量采集、\code{shift score} 接入最终异常分数、平滑、阈值化；
    \item \code{mra.py:563--636}：神经网络联合训练入口，用来说明哪些参数通过反向传播更新；
    \item \code{mra.py:954--1029}：训练后的补全序列如何被重新窗口化，并成为 \code{shift score} 的输入。
\end{itemize}

其它模块只在必要处说明。比如双专家插补器和 Transformer 检测器会影响 \code{shift score} 的输入，但它们自身的完整结构不是本文重点。

\subsection{先给结论}

结合代码阅读和默认数据形状探针，可以得到以下核心结论：
\begin{itemize}
    \item \code{shift score} 的代码入口是 \code{compute_distribution_shift_scores}，位于 \code{mra.py:830--843}；它使用的特征入口是 \code{build_window_feature_matrix}，位于 \code{mra.py:815--827}。
    \item \code{shift score} 的主流程入口在 \code{main} 中，位于 \code{mra.py:1012--1029}。代码先把训练集和测试集的补全序列重新切成评分窗口，再计算训练窗口和测试窗口的分布偏移分数。
    \item \code{shift score} 本身没有神经网络参数，也没有 \code{loss.backward()} 和 \code{optimizer.step()}。它的“训练”是统计意义上的拟合：用训练集窗口特征估计均值 \(\bm{\mu}_f\) 和标准差 \(\bm{\sigma}_f\)。
    \item 神经网络参数真正更新的位置是 \code{mra.py:563--636} 的 \code{train_model}。这里使用 AdamW 同时更新图专家、频域专家、门控网络和 Transformer 检测器。
    \item \code{shift score} 并不直接使用原始含缺失值的窗口，而是使用训练后模型生成的稠密补全序列 \code{train_complete_scaled} 和 \code{test_complete_scaled}。
    \item 默认数据有 \(18\) 个变量，窗口长度为 \(50\)。训练评分窗口为 \(4001 \times 50 \times 18\)，测试评分窗口为 \(4000 \times 50 \times 18\)。\code{shift score} 每个窗口提取 \(108\) 维统计特征。
    \item \code{shift score} 先计算每个窗口相对训练特征均值的标准化距离，再进入 \code{combine_scores}。在最终分数里，它与 \code{detector_score} 和 \code{disagreement_score} 合成。
    \item \code{shift score} 的训练集分数不是异常标签，而是校准基准。训练集分数用于估计该分数通道自己的均值、标准差，并参与训练分数阈值的确定。
    \item 当前实现中，\code{shift score} 的特征均值和标准差是函数内局部变量，没有保存到 \code{model.pt}。如果以后要把训练和测试拆成两个独立进程，需要额外保存这两个统计量。
\end{itemize}

\section{代码入口、训练入口与总体调用链}

\subsection{入口总表}

{\small
\begin{longtable}{@{}>{\raggedright\arraybackslash}p{0.23\linewidth} >{\raggedright\arraybackslash}p{0.23\linewidth} >{\raggedright\arraybackslash}p{0.46\linewidth}@{}}
\toprule
入口 & 代码位置 & 作用 \\
\midrule
\endhead
命令行权重入口 & \codecell{mra.py:104--108} & 定义 \code{--score-shift-weight}，控制 \code{shift score} 在最终异常分数中的权重，默认值为 \(1.0\)。 \\
训练窗口入口 & \codecell{mra.py:875--898} & 构造插补训练窗口、训练评分窗口和测试评分窗口。 \\
神经网络训练入口 & \codecell{mra.py:927--941} 调用 \codecell{mra.py:563--636} & 训练双专家插补器和 Transformer 检测器，更新神经网络权重。 \\
补全序列入口 & \codecell{mra.py:954--973} & 用训练后的模型生成 \code{train_complete_scaled} 和 \code{test_complete_scaled}。 \\
基础统计分数入口 & \codecell{mra.py:993--1010} 调用 \codecell{mra.py:680--737} & 计算 \code{detector_score}、\code{disagreement_score} 和 \code{gate_entropy}。 \\
\code{shift score} 窗口入口 & \codecell{mra.py:1012--1023} & 把补全后的训练集和测试集重新切成无缺失评分窗口。 \\
\code{shift score} 计算入口 & \codecell{mra.py:1024--1029} 调用 \codecell{mra.py:830--843} & 用训练集补全窗口拟合分布统计量，并给训练/测试窗口计算偏移分数。 \\
特征构造入口 & \codecell{mra.py:815--827} & 从每个窗口中提取均值、标准差、一阶变化和 FFT 幅值特征。 \\
分数合成入口 & \codecell{mra.py:1031--1036} 调用 \codecell{mra.py:752--773} & 把三类分数统一到训练集尺度后合成最终原始分数。 \\
平滑和阈值入口 & \codecell{mra.py:1037--1055} & 对分数做 EWAF 平滑，用训练分数确定阈值，再生成测试预测。 \\
\bottomrule
\end{longtable}
}

\subsection{两类“训练入口”必须分开}

这份代码里有两种容易混淆的“训练”。

\paragraph{第一类：神经网络训练。}
这是真正的反向传播训练，入口是 \code{train_model}。每个 batch 中，代码随机遮挡一部分原本可观测的位置，令模型去恢复这些位置。随后计算五项损失，加权求和，执行反向传播和 AdamW 更新。对应代码为：
\begin{equation}
\begin{aligned}
    \texttt{outputs = model(...)}
    &\rightarrow
    \texttt{losses = compute\_losses(...)} \\
    &\rightarrow
    \texttt{backward()}
    \rightarrow
    \texttt{optimizer.step()}.
\end{aligned}
\end{equation}

\paragraph{第二类：\code{shift score} 的统计拟合。}
\code{shift score} 不执行反向传播。它的拟合过程只有三步：
\begin{enumerate}
    \item 从训练集补全窗口中提取特征矩阵 \(\bm{F}^{tr}\)；
    \item 计算训练特征均值 \(\bm{\mu}_f\) 和标准差 \(\bm{\sigma}_f\)；
    \item 用 \(\bm{\mu}_f,\bm{\sigma}_f\) 标准化训练窗口和测试窗口，并计算每个窗口的标准化距离。
\end{enumerate}

因此，若把 \code{shift score} 称为“训练过的模块”，这里的“训练”应理解为“用训练集估计统计基准”，而不是“用梯度下降学习参数”。

\subsection{完整调用链}

从原始数据到最终测试预测，和 \code{shift score} 相关的实际链路如下：
\begin{equation}
\begin{aligned}
    &\text{读取 CSV}
    \rightarrow
    \text{观测值标准化}
    \rightarrow
    \text{构造训练插补窗口}
    \rightarrow
    \text{训练联合模型} \\
    &\rightarrow
    \text{补全训练/测试全序列}
    \rightarrow
    \text{把补全序列切成评分窗口}
    \rightarrow
    \text{提取窗口特征} \\
    &\rightarrow
    \text{用训练特征拟合均值和标准差}
    \rightarrow
    \text{计算训练/测试 shift score}
    \rightarrow
    \text{合成最终异常分数}.
\end{aligned}
\end{equation}

这条链路说明一个关键事实：\code{shift score} 虽然没有梯度训练，但它依赖前面已经训练好的插补模型。插补模型变了，补全序列会变，窗口特征会变，\code{shift score} 也会随之改变。

\section{数据、窗口与张量形状}

\subsection{默认数据形状}

默认命令读取：
\begin{itemize}
    \item \code{data/train/train_1.csv}；
    \item \code{data/test/test_C5_1.csv}。
\end{itemize}

当前仓库数据的实际形状为：
\begin{align}
    \bm{X}^{tr}_{raw} &\in \mathbb{R}^{4100 \times 18}, \\
    \bm{X}^{te}_{raw} &\in \mathbb{R}^{4100 \times 18}, \\
    \bm{M}^{tr}_{s},\bm{M}^{te}_{s} &\in \{0,1\}^{4100 \times 18}.
\end{align}

其中 \(\bm{M}_s\) 是结构性缺失掩码。若 \(M_{s,t,n}=1\)，表示第 \(t\) 个时间点、第 \(n\) 个变量在原始 CSV 中是 \code{NaN}。

\subsection{标准化器对 \code{shift score} 的影响}

\code{ObservedStandardScaler} 在 \code{mra.py:865--867} 被调用。它只用训练集观测值拟合每个变量的均值和标准差：
\begin{equation}
    \mu_n =
    \frac{1}{|\Omega_n^{tr}|}
    \sum_{t\in\Omega_n^{tr}} X^{tr}_{raw,t,n},
    \qquad
    \sigma_n =
    \max\left(
    \sqrt{\frac{1}{|\Omega_n^{tr}|}
    \sum_{t\in\Omega_n^{tr}}(X^{tr}_{raw,t,n}-\mu_n)^2},
    10^{-6}
    \right),
\end{equation}
其中 \(\Omega_n^{tr}\) 表示训练集中第 \(n\) 个变量的可观测位置集合。

训练集和测试集随后都使用同一组 \(\mu_n,\sigma_n\) 变换：
\begin{equation}
    X_{t,n} =
    \begin{cases}
        \dfrac{X_{raw,t,n}-\mu_n}{\sigma_n}, & M_{s,t,n}=0, \\[6pt]
        0, & M_{s,t,n}=1.
    \end{cases}
\end{equation}

这对 \code{shift score} 很重要，因为 \code{shift score} 使用的是 \code{train_complete_scaled} 和 \code{test_complete_scaled}。也就是说，分布偏移是在训练集标准化坐标系里衡量的。

\subsection{窗口形状}

默认参数为 \(\code{seq-len}=50\)、\(\code{stride}=1\)。代码会构造两类窗口。

\paragraph{插补窗口。}
插补窗口由 \code{build_windows} 构造，用于神经网络训练和全序列补全。它允许在序列开头重复首行补齐，因此窗口数量等于原始时间步数：
\begin{align}
    \bm{W}^{tr}_{imp} &\in \mathbb{R}^{4100 \times 50 \times 18}, \\
    \bm{W}^{te}_{imp} &\in \mathbb{R}^{4100 \times 50 \times 18}.
\end{align}

\paragraph{评分窗口。}
评分窗口由 \code{build_standard_windows} 和 \code{build_prompt_test_windows} 构造，用于计算训练分数和测试分数。默认数据下：
\begin{align}
    \bm{W}^{tr}_{eval} &\in \mathbb{R}^{4001 \times 50 \times 18}, \\
    \bm{W}^{te}_{eval} &\in \mathbb{R}^{4000 \times 50 \times 18}.
\end{align}

测试标签由 \code{build_prompt_test_windows} 生成，共 \(4000\) 个窗口，前 \(2000\) 个为正常，后 \(2000\) 个为异常。

\subsection{shift score 使用的窗口不是原始评分窗口}

\code{detector_score} 和 \code{disagreement_score} 使用的是 \(\bm{W}^{tr}_{eval}\) 与 \(\bm{W}^{te}_{eval}\)。但 \code{shift score} 另走一条路径：
\begin{enumerate}
    \item 先用训练后的模型生成补全序列：
    \begin{equation}
        \bm{C}^{tr},\bm{C}^{te}\in\mathbb{R}^{4100\times18}.
    \end{equation}
    \item 再把补全序列重新切成窗口：
    \begin{align}
        \bm{V}^{tr} &= \operatorname{Window}(\bm{C}^{tr})
        \in \mathbb{R}^{4001 \times 50 \times 18}, \\
        \bm{V}^{te} &= \operatorname{Window}(\bm{C}^{te})
        \in \mathbb{R}^{4000 \times 50 \times 18}.
    \end{align}
    \item 切窗口时传入的掩码是全零矩阵，即代码中的 \code{np.zeros_like(train_mask)} 和 \code{np.zeros_like(test_mask)}。
\end{enumerate}

这说明 \code{shift score} 衡量的是“模型补全后的稠密窗口是否像训练集补全窗口”，而不是“原始缺失模式是否像训练集”。这样做可以避免多采样率缺失模式直接主导分数，把注意力放在补全后的数据分布上。

\section{联合模型训练流程与参数更新}

\subsection{训练 batch 内部的数据流}

在 \code{train_model} 中，一个训练 batch 的张量记为：
\begin{equation}
    \bm{X}_{true}\in\mathbb{R}^{B\times T\times N},
    \qquad
    \bm{M}_s\in\{0,1\}^{B\times T\times N}.
\end{equation}
默认 \(T=50\)、\(N=18\)，\(B\) 由 \code{batch-size} 决定，默认最多为 \(64\)。

训练时先从真实观测位置随机抽出一部分作为自监督目标：
\begin{equation}
    \bm{M}_h = \operatorname{SampleHoldout}(\bm{M}_s),
    \qquad
    \bm{M} = \max(\bm{M}_s,\bm{M}_h).
\end{equation}
其中 \(\bm{M}_h\) 是训练期额外遮挡掩码，\(\bm{M}\) 是模型实际看到的缺失掩码。

输入窗口为：
\begin{equation}
    \bm{X}_{input} =
    \bm{X}_{true}\odot(1-\bm{M}).
\end{equation}
随后模型内部执行：
\begin{align}
    \bm{X}_{seed} &= \operatorname{PreImpute}(\bm{X}_{input},\bm{M}), \\
    \bm{X}_{gcn} &= f_{gcn}(\bm{X}_{seed},\bm{M}), \\
    \bm{X}_{freq} &= f_{freq}(\bm{X}_{seed}), \\
    \bm{\Pi} &= \operatorname{softmax}(g(\bm{X}_{seed},\bm{M},\bm{X}_{gcn},\bm{X}_{freq},\bm{rate\_id})), \\
    \bm{X}_{imp} &= \Pi_{0}\odot\bm{X}_{gcn}+\Pi_{1}\odot\bm{X}_{freq}, \\
    \bm{X}_{complete} &= \bm{M}\odot\bm{X}_{imp}+(1-\bm{M})\odot\bm{X}_{input}, \\
    \widehat{\bm{X}}_{det} &= f_{det}(\bm{X}_{complete},1-\bm{M}_s,\bm{rate\_id},\bm{stride}).
\end{align}

\subsection{主要输出张量形状}

{\small
\begin{longtable}{@{}>{\raggedright\arraybackslash}p{0.25\linewidth} >{\raggedright\arraybackslash}p{0.24\linewidth} >{\raggedright\arraybackslash}p{0.43\linewidth}@{}}
\toprule
张量 & 形状 & 含义 \\
\midrule
\endhead
\(\bm{X}_{true}\) & \(B\times T\times N\) & 训练 batch 的标准化真实窗口。 \\
\(\bm{M}_s\) & \(B\times T\times N\) & 原始结构性缺失掩码。 \\
\(\bm{M}_h\) & \(B\times T\times N\) & 训练期随机遮挡出的 holdout 位置。 \\
\(\bm{M}\) & \(B\times T\times N\) & 模型实际看到的缺失掩码，等于结构性缺失和 holdout 的并集。 \\
\(\bm{X}_{seed}\) & \(B\times T\times N\) & 双向填充后的时域种子。 \\
\(\bm{X}_{gcn}\) & \(B\times T\times N\) & 图专家重构结果。 \\
\(\bm{X}_{freq}\) & \(B\times T\times N\) & 频域专家重构结果。 \\
\(\bm{\Pi}\) & \(B\times T\times N\times2\) & 每个位置对 GCN 专家和频域专家的门控权重。 \\
\(\bm{A}\) & \(B\times N\times N\) & 图专家生成的窗口级邻接矩阵。 \\
\(\bm{X}_{complete}\) & \(B\times T\times N\) & 观测位置保留原值、缺失位置采用门控补全后的完整窗口。 \\
\(\widehat{\bm{X}}_{det}\) & \(B\times T\times N\) & Transformer 检测器对完整窗口的重构。 \\
\bottomrule
\end{longtable}
}

\subsection{损失函数和参数更新}

\code{compute_losses} 定义五项损失：
\begin{align}
    \mathcal{L}_{fusion}
    &=
    \frac{\sum ((\bm{X}_{imp}-\bm{X}_{true})\odot\bm{M}_h)^2}
    {\sum \bm{M}_h}, \\
    \mathcal{L}_{gcn}
    &=
    \frac{\sum ((\bm{X}_{gcn}-\bm{X}_{true})\odot\bm{M}_h)^2}
    {\sum \bm{M}_h}, \\
    \mathcal{L}_{freq}
    &=
    \frac{\sum ((\bm{X}_{freq}-\bm{X}_{true})\odot\bm{M}_h)^2}
    {\sum \bm{M}_h}, \\
    \mathcal{L}_{graph}
    &=
    \operatorname{mean}\left((\operatorname{diag}(\bm{A})-\tau)^2\right), \\
    \mathcal{L}_{det}
    &=
    \frac{\sum ((\widehat{\bm{X}}_{det}-\sg(\bm{X}_{complete}))\odot(1-\bm{M}_s))^2}
    {\sum (1-\bm{M}_s)}.
\end{align}

其中 \(\sg(\cdot)\) 表示 \code{detach()}。它让检测器目标值不再反向传播，但 \(\bm{X}_{complete}\) 作为检测器输入仍然会影响 \(\widehat{\bm{X}}_{det}\)，所以 \(\mathcal{L}_{det}\) 仍可能反向更新插补器。

总损失为：
\begin{equation}
    \mathcal{L}
    =
    1.0\mathcal{L}_{fusion}
    +0.4\mathcal{L}_{gcn}
    +0.4\mathcal{L}_{freq}
    +0.5\mathcal{L}_{det}
    +0.1\mathcal{L}_{graph}.
\end{equation}

参数更新流程为：
\begin{enumerate}
    \item \code{optimizer.zero_grad()} 清空旧梯度；
    \item \code{losses["total_loss"].backward()} 计算梯度；
    \item \code{clip_grad_norm_(model.parameters(), max_norm=1.0)} 做梯度裁剪；
    \item \code{optimizer.step()} 用 AdamW 更新神经网络参数。
\end{enumerate}

\code{shift score} 不在这四步里。它没有梯度，也不会改变神经网络权重；它只在训练结束后读取补全结果并估计统计基准。

\section{shift score 的数学原理}

\subsection{它想解决什么问题}

\code{detector_score} 衡量的是检测器重构误差，\code{disagreement_score} 衡量的是两个插补专家之间的分歧。这两类分数都很有用，但仍有一个缺口：如果测试窗口整体分布已经偏离训练集，但模型仍能给出平滑补全，重构误差和专家分歧未必会很大。

\code{shift score} 用来补这个缺口。它不问“模型重构错了多少”，而问：
\begin{quote}
    这个补全后的窗口，在均值、波动、变化速度和频域能量上，还像不像训练集窗口？
\end{quote}

如果一个窗口在这些统计特征上离训练集常态很远，它的 \code{shift score} 就会变大。

\subsection{窗口特征构造}

设一个补全窗口为：
\begin{equation}
    \bm{V}_i\in\mathbb{R}^{T\times N},
\end{equation}
其中 \(i\) 表示第 \(i\) 个窗口。默认 \(T=50\)、\(N=18\)。

\code{build_window_feature_matrix} 提取四类特征。

\paragraph{第一类：变量组均值。}
代码固定使用四个列切片：
\begin{equation}
    G_1=\{1,\ldots,6\},\quad
    G_2=\{7,\ldots,12\},\quad
    G_3=\{13,\ldots,18\},\quad
    G_4=\{1,\ldots,18\}.
\end{equation}
对每个组内变量，计算窗口时间均值：
\begin{equation}
    m_{i,n}=\frac{1}{T}\sum_{t=1}^{T} V_{i,t,n}.
\end{equation}
这类特征描述窗口整体水平是否偏高或偏低。

\paragraph{第二类：变量组标准差。}
对同样四个列切片，计算时间标准差：
\begin{equation}
    s_{i,n}=
    \sqrt{\frac{1}{T}\sum_{t=1}^{T}(V_{i,t,n}-m_{i,n})^2}.
\end{equation}
这类特征描述窗口内部波动是否变强或变弱。

\paragraph{第三类：一阶变化幅度。}
代码计算相邻时间点差分的平均绝对值：
\begin{equation}
    d_{i,n}=
    \frac{1}{T-1}\sum_{t=2}^{T}|V_{i,t,n}-V_{i,t-1,n}|.
\end{equation}
它描述曲线在窗口内变化得快不快。如果某个变量突然抖动、跳变或振荡增强，这个特征会变大。

\paragraph{第四类：频域幅值。}
代码沿时间维做实数快速傅里叶变换：
\begin{equation}
    \widehat{V}_{i,k,n}
    =
    \sum_{t=0}^{T-1}V_{i,t,n}
    \exp\left(-\frac{2\pi \mathrm{i}kt}{T}\right).
\end{equation}
随后取第 \(1\) 到第 \(5\) 个非零频点的幅值均值：
\begin{equation}
    a_{i,n}
    =
    \frac{1}{5}\sum_{k=1}^{5}|\widehat{V}_{i,k,n}|.
\end{equation}
这里跳过第 \(0\) 个频点，是因为第 \(0\) 个频点主要表示窗口平均水平，而平均水平已经由前面的均值特征表达。非零频点更能表达周期性和振荡强度。

\subsection{特征维度拆解}

默认 \(N=18\) 时，每个窗口得到 \(108\) 维特征。拆解如下。

{\small
\begin{longtable}{@{}>{\raggedright\arraybackslash}p{0.22\linewidth} >{\raggedright\arraybackslash}p{0.34\linewidth} >{\raggedright\arraybackslash}p{0.16\linewidth} >{\raggedright\arraybackslash}p{0.20\linewidth}@{}}
\toprule
特征来源 & 代码对应 & 维度 & 通俗含义 \\
\midrule
\endhead
第 1--6 列均值和标准差 & \codecell{slice(0, 6)} & \(6+6\) & 高频变量的水平和波动。 \\
第 7--12 列均值和标准差 & \codecell{slice(6, 12)} & \(6+6\) & 中频变量的水平和波动。 \\
第 13--18 列均值和标准差 & \codecell{slice(12, 18)} & \(6+6\) & 低频变量的水平和波动。 \\
第 1--18 列均值和标准差 & \codecell{slice(0, 18)} & \(18+18\) & 全变量整体水平和波动。 \\
相邻时间差分平均绝对值 & \codecell{abs(diff(windows)).mean(axis=1)} & \(18\) & 曲线变化快慢。 \\
FFT 第 1--5 个非零频点平均幅值 & \codecell{abs(rfft(windows))[:, 1:6, :].mean(axis=1)} & \(18\) & 低频到中低频振荡强度。 \\
\midrule
合计 & -- & \(108\) & 每个窗口的统计画像。 \\
\bottomrule
\end{longtable}
}

注意，这里的前三个变量组和总变量组是按当前 \(18\) 列数据写死的。它体现了当前数据的三类采样率结构：前 \(6\) 列高频，中间 \(6\) 列中频，最后 \(6\) 列低频。若未来数据列数或分组方式变化，这部分代码需要同步改造。

\subsection{从窗口特征到分布偏移分数}

设训练补全窗口数为 \(K_{tr}\)，测试补全窗口数为 \(K_{te}\)。默认情况下：
\begin{equation}
    K_{tr}=4001,\qquad K_{te}=4000.
\end{equation}

特征矩阵为：
\begin{equation}
    \bm{F}^{tr}\in\mathbb{R}^{K_{tr}\times P},
    \qquad
    \bm{F}^{te}\in\mathbb{R}^{K_{te}\times P},
    \qquad
    P=108.
\end{equation}

\code{compute_distribution_shift_scores} 先用训练特征估计每一维特征的均值和标准差：
\begin{align}
    \mu_{f,j}
    &=
    \frac{1}{K_{tr}}\sum_{i=1}^{K_{tr}}F^{tr}_{i,j}, \\
    \sigma_{f,j}
    &=
    \max\left(
    \sqrt{
    \frac{1}{K_{tr}}\sum_{i=1}^{K_{tr}}(F^{tr}_{i,j}-\mu_{f,j})^2
    },
    10^{-6}
    \right).
\end{align}

然后训练窗口和测试窗口都用同一组训练统计量做标准化：
\begin{align}
    Z^{tr}_{i,j}
    &=
    \frac{F^{tr}_{i,j}-\mu_{f,j}}{\sigma_{f,j}}, \\
    Z^{te}_{i,j}
    &=
    \frac{F^{te}_{i,j}-\mu_{f,j}}{\sigma_{f,j}}.
\end{align}

最后对每个窗口计算所有特征维度的均方根距离：
\begin{align}
    q^{tr}_i
    &=
    \sqrt{\frac{1}{P}\sum_{j=1}^{P}(Z^{tr}_{i,j})^2}, \\
    q^{te}_i
    &=
    \sqrt{\frac{1}{P}\sum_{j=1}^{P}(Z^{te}_{i,j})^2}.
\end{align}

这里的 \(q_i\) 就是原始 \code{shift_score}。

\subsection{通俗解释}

可以把每个窗口想成一张“体检表”。这张表有 \(108\) 个指标，包括平均值、波动大小、变化速度和频域振荡强度。训练集给出了正常情况下每个指标的大致范围。测试时，如果某个窗口的很多指标都偏离训练集正常范围，那么它的 \code{shift score} 会升高。

公式中的 \(Z_{i,j}\) 是常见的 z 分数。它表示第 \(i\) 个窗口第 \(j\) 个特征离训练均值有几个训练标准差。例如：
\begin{itemize}
    \item \(Z_{i,j}=0\)：这个特征刚好等于训练平均水平；
    \item \(Z_{i,j}=1\)：这个特征比训练均值高一个训练标准差；
    \item \(Z_{i,j}=-2\)：这个特征比训练均值低两个训练标准差。
\end{itemize}

\code{shift score} 对所有 \(Z_{i,j}\) 平方后取平均，再开方。因此它不关心偏高还是偏低，只关心偏离程度。一个窗口只要在很多特征上都偏离训练集，就会得到较大的分数。

\subsection{与马氏距离的关系}

\code{shift score} 可以看成一种简化的标准化距离。它和马氏距离的共同点是：都先用训练集估计正常分布，再衡量新样本离正常分布有多远。

差别是：当前实现只除以每一维自己的标准差，没有使用完整协方差矩阵。也就是说，它默认这些统计特征之间相互独立，或者至少不显式建模特征间相关性。因此它更简单、更稳定，也更不容易在样本数有限时出现协方差矩阵求逆问题；代价是不能表达“两个特征单独看正常，但组合起来不正常”的情况。

\section{shift score 的输入、输出与内部数据流}

\subsection{函数级输入输出}

{\small
\begin{longtable}{@{}>{\raggedright\arraybackslash}p{0.22\linewidth} >{\raggedright\arraybackslash}p{0.30\linewidth} >{\raggedright\arraybackslash}p{0.40\linewidth}@{}}
\toprule
函数 & 输入/输出形状 & 说明 \\
\midrule
\endhead
\makecell[l]{\ttfamily\small build\_window\\\ttfamily\small feature\_matrix} & 输入：\(K\times T\times N\)。输出：\(K\times P\)。默认 \(P=108\)。 & 每个窗口变成一行统计特征。时间维被压缩掉。 \\
\makecell[l]{\ttfamily\small compute\_distribution\\\ttfamily\small shift\_scores} & 输入：\(\bm{V}^{tr}\in\mathbb{R}^{K_{tr}\times T\times N}\)，\(\bm{V}^{eval}\in\mathbb{R}^{K_{eval}\times T\times N}\)。输出：\(\bm{q}^{tr}\in\mathbb{R}^{K_{tr}}\)，\(\bm{q}^{eval}\in\mathbb{R}^{K_{eval}}\)。 & 用训练窗口拟合特征均值和标准差，再同时给训练集和评估集打分。 \\
\makecell[l]{\ttfamily\small normalize\\\ttfamily\small scores} & 输入：两个一维分数数组。输出：两个同形状 z 分数数组。 & 再次用训练分数均值和标准差统一分数通道尺度。 \\
\makecell[l]{\ttfamily\small combine\\\ttfamily\small scores} & 输入：包含三类分数的一维数组字典。输出：训练原始总分和评估原始总分。 & 把检测器误差、专家分歧和分布偏移合成。 \\
\bottomrule
\end{longtable}
}

\subsection{主流程中的实际形状}

默认数据和默认窗口参数下，\code{shift score} 的实际形状如下：
\begin{align}
    \code{train_complete_scaled} &:\quad 4100\times18, \\
    \code{test_complete_scaled} &:\quad 4100\times18, \\
    \code{train_complete_eval_windows} &:\quad 4001\times50\times18, \\
    \code{test_complete_eval_windows} &:\quad 4000\times50\times18, \\
    \code{train_features} &:\quad 4001\times108, \\
    \code{test_features} &:\quad 4000\times108, \\
    \code{feature_mean},\code{feature_std} &:\quad 108, \\
    \code{train_shift_scores} &:\quad 4001, \\
    \code{test_shift_scores} &:\quad 4000.
\end{align}

\subsection{为什么要先补全再计算}

原始多采样率数据中有大量 \code{NaN}。如果直接对原始窗口提取均值、标准差和 FFT 特征，会遇到两个问题：
\begin{enumerate}
    \item \code{NaN} 无法直接参与普通的 \code{mean/std/fft} 计算；
    \item 如果简单把缺失值填零，分数可能主要反映采样率缺失模式，而不是真实数据分布变化。
\end{enumerate}

当前实现先用训练好的双专家模型补全序列，再计算 \code{shift score}。这样得到的分数更接近“补全后的过程状态是否偏离训练集”，而不是“缺失模式是否偏离训练集”。

\section{异常分数合成与测试集作用}

\subsection{三类分数}

\code{mra.py} 中用于最终合成的三类分数为：
\begin{enumerate}
    \item \code{detector_score}：检测器重构误差；
    \item \code{disagreement_score}：GCN 专家和频域专家输出的平均绝对差；
    \item \code{shift_score}：补全窗口相对训练补全窗口的分布偏移距离。
\end{enumerate}

\code{gate_entropy} 也会被计算并写入 \code{test_predictions.csv}，但当前代码没有把它并入最终异常分数。

\subsection{检测器分数}

在 \code{collect_window_statistics} 中，检测器分数为：
\begin{equation}
    r_i =
    \frac{
    \sum_{t,n}
    \left((\widehat{X}_{det,i,t,n}-X_{complete,i,t,n})O_{i,t,n}\right)^2
    }{
    \sum_{t,n}O_{i,t,n}
    },
\end{equation}
其中 \(O=1-M_s\) 是结构性观测掩码。它只在原始可观测位置上计算误差。

\subsection{专家分歧分数}

专家分歧分数为：
\begin{equation}
    g_i =
    \frac{1}{TN}\sum_{t,n}|X_{gcn,i,t,n}-X_{freq,i,t,n}|.
\end{equation}
它表示两个插补专家对同一个窗口的看法是否一致。如果两个专家差距很大，说明模型对这个窗口不够确定。

\subsection{shift score 的合成方式}

进入 \code{combine_scores} 后，三个分数通道都会用训练集的通道均值和标准差做归一化：
\begin{align}
    \widetilde{r}_i &=
    \frac{r_i-\mu_r^{tr}}{\sigma_r^{tr}}, \\
    \widetilde{g}_i &=
    \frac{g_i-\mu_g^{tr}}{\sigma_g^{tr}}, \\
    \widetilde{q}_i &=
    \frac{q_i-\mu_q^{tr}}{\sigma_q^{tr}}.
\end{align}

最终原始异常分数为：
\begin{equation}
    S_i =
    |\widetilde{r}_i|
    +
    \lambda_g\widetilde{g}_i
    +
    \lambda_q\widetilde{q}_i.
\end{equation}
默认：
\begin{equation}
    \lambda_g=0.25,
    \qquad
    \lambda_q=1.0.
\end{equation}

这里有一个实现细节需要注意：代码只对 \(\widetilde{r}_i\) 取绝对值，没有对 \(\widetilde{g}_i\) 和 \(\widetilde{q}_i\) 取绝对值。因此，如果某个测试窗口的 \code{shift score} 低于训练集平均水平，\(\widetilde{q}_i\) 会是负数，会降低最终分数。这并不是数学错误，而是当前实现的设计选择：高于训练常态的分布偏移会加分，低于训练常态的偏移会减分。

\subsection{平滑、阈值与最终预测}

合成后的训练分数和测试分数会经过 EWAF 平滑：
\begin{equation}
    \bar{S}_i =
    \alpha S_i + (1-\alpha)\bar{S}_{i-1},
    \qquad
    \alpha=\code{ewaf-alpha}.
\end{equation}
默认 \(\alpha=0.3\)。

阈值只由训练分数决定：
\begin{equation}
    \theta =
    \max\left(
    \operatorname{mean}(\bar{\bm{S}}^{tr})
    + k\operatorname{std}(\bar{\bm{S}}^{tr}),
    \operatorname{Quantile}_{p}(\bar{\bm{S}}^{tr})
    \right),
\end{equation}
默认 \(k=2.0\)、\(p=0.95\)。

测试窗口先按阈值得到初始预测：
\begin{equation}
    \hat{y}_i = \mathbb{I}(\bar{S}^{te}_i\ge \theta).
\end{equation}
然后 \code{apply_min_anomaly_duration} 会删除短于 \code{min-anomaly-duration} 的异常片段，默认最短持续长度为 \(50\) 个窗口。

\subsection{训练集统计量在测试集中的作用}

训练集估计出的量不是只服务训练，它们会贯穿整个测试流程。

{\small
\begin{longtable}{@{}>{\raggedright\arraybackslash}p{0.22\linewidth} >{\raggedright\arraybackslash}p{0.25\linewidth} >{\raggedright\arraybackslash}p{0.45\linewidth}@{}}
\toprule
训练集产生的量 & 代码位置 & 在测试集中的作用 \\
\midrule
\endhead
标准化均值和标准差 \(\mu_n,\sigma_n\) & \codecell{mra.py:865--867} & 测试集用训练集尺度标准化，保证训练和测试处在同一个数值坐标系。 \\
距离先验矩阵 & \codecell{mra.py:868--873} & 若未提供外部物理距离，图先验由训练集数据自动计算，并影响测试时 GCN 专家输出。 \\
\code{rate_id} 和 \code{stride} & \codecell{mra.py:900} & 测试时仍使用训练集推断出的采样率类别和步长，供门控网络和检测器编码使用。 \\
神经网络权重 \(\bm{\theta}\) & \codecell{mra.py:583, 618--621} & 测试补全、检测器重构、专家分歧和 \code{shift score} 的输入都依赖训练后的权重。 \\
\code{shift score} 特征均值 \(\bm{\mu}_f\) & \codecell{mra.py:836} & 测试窗口特征要减去训练均值，衡量离训练常态的偏离。 \\
\code{shift score} 特征标准差 \(\bm{\sigma}_f\) & \codecell{mra.py:837} & 测试窗口特征要除以训练标准差，使不同特征维度可比较。 \\
各分数通道的训练均值和标准差 & \codecell{mra.py:740--749} & 测试分数通道用训练集尺度归一化，再进入最终分数合成。 \\
异常阈值 \(\theta\) & \codecell{mra.py:1045--1050} & 测试分数是否异常完全由训练分数确定的阈值判断。 \\
\bottomrule
\end{longtable}
}

\section{代码实现逐段说明}

\subsection{\texorpdfstring{\texttt{build\_window\_feature\_matrix}}{build window feature matrix}}

核心代码逻辑可以概括为：
\begin{enumerate}
    \item 初始化空列表 \code{features}；
    \item 对四个列切片分别取子窗口；
    \item 对每个子窗口沿时间维计算均值和标准差；
    \item 对全变量窗口计算相邻差分平均绝对值；
    \item 对全变量窗口沿时间维做 \code{rfft}，取第 \(1\) 到第 \(5\) 个频点的平均幅值；
    \item 沿特征维拼接所有结果，返回 \code{float32} 矩阵。
\end{enumerate}

从张量角度看，若输入为：
\begin{equation}
    \bm{V}\in\mathbb{R}^{K\times50\times18},
\end{equation}
那么：
\begin{itemize}
    \item \code{block.mean(axis=1)} 把时间维 \(50\) 压缩掉；
    \item \code{block.std(axis=1)} 同样把时间维压缩掉；
    \item \code{np.diff(windows, axis=1)} 把时间维从 \(50\) 变为 \(49\)，再 \code{mean(axis=1)} 压缩为 \(K\times18\)；
    \item \code{np.fft.rfft(windows, axis=1)} 在长度 \(50\) 上产生 \(26\) 个频点，切片 \code{[:, 1:6, :]} 取 \(5\) 个非零频点，再对这 \(5\) 个频点求均值，得到 \(K\times18\)。
\end{itemize}

\subsection{\texorpdfstring{\texttt{compute\_distribution\_shift\_scores}}{compute distribution shift scores}}

该函数的输入是训练窗口和评估窗口。这里的 \code{eval_windows} 可以是测试窗口，也可以是其它需要评估的窗口。

函数内部有四个关键变量：
\begin{align}
    \code{train_features} &\leftrightarrow \bm{F}^{tr}, \\
    \code{eval_features} &\leftrightarrow \bm{F}^{eval}, \\
    \code{feature_mean} &\leftrightarrow \bm{\mu}_f, \\
    \code{feature_std} &\leftrightarrow \bm{\sigma}_f.
\end{align}

它返回两个一维数组：
\begin{align}
    \code{train_score} &\in\mathbb{R}^{K_{tr}}, \\
    \code{eval_score} &\in\mathbb{R}^{K_{eval}}.
\end{align}

这两个数组都使用同一组 \(\bm{\mu}_f,\bm{\sigma}_f\)。这保证训练分数和测试分数可比较。

\subsection{\texorpdfstring{\texttt{main} 中的接入位置}{main 中的接入位置}}

\code{main} 中与 \code{shift score} 直接相关的关键语句按顺序是：
\begin{enumerate}
    \item \code{reconstruct_full_sequence(...)} 生成训练集和测试集补全序列；
    \item \code{build_standard_windows(train_complete_scaled, zeros)} 生成训练补全评分窗口；
    \item \code{build_prompt_test_windows(test_complete_scaled, zeros)} 生成测试补全评分窗口；
    \item \code{compute_distribution_shift_scores(...)} 计算训练和测试 \code{shift score}；
    \item 把结果写入 \code{train_stats["shift_score"]} 和 \code{test_stats["shift_score"]}；
    \item \code{combine_scores(..., shift_weight=args.score_shift_weight)} 将其并入最终分数。
\end{enumerate}

这里传入全零掩码的原因是：补全序列已经没有结构性缺失，重新窗口化时不应再把任何位置当作缺失位置。

\section{设计思路、优点与局限}

\subsection{设计思路}

\code{shift score} 的设计可以概括为一句话：
\begin{quote}
    用训练集补全窗口建立一个“正常窗口统计画像”，再看测试补全窗口偏离这张画像多少。
\end{quote}

它刻意选择了容易解释的统计量，而不是再训练一个复杂分类器：
\begin{itemize}
    \item 均值特征看窗口整体水平；
    \item 标准差特征看窗口内部波动；
    \item 差分特征看相邻时间变化速度；
    \item FFT 幅值特征看周期性或振荡强度。
\end{itemize}

这些特征都能用简单语言解释，因此适合放进异常检测报告里说明。

\subsection{为什么和重构误差互补}

重构误差依赖模型是否能重建当前窗口。如果模型泛化能力很强，它可能把异常窗口也重构得不错，导致 \code{detector_score} 不高。专家分歧也类似，如果两个专家都朝同一个异常模式输出，分歧也可能不明显。

\code{shift score} 不直接看模型是否重构成功，而是看补全后的窗口统计形态是否偏离训练集。因此它能补充另一种异常信号：过程整体分布变化。

\subsection{当前实现的优点}

\begin{itemize}
    \item 计算简单，不需要额外神经网络；
    \item 使用训练集统计量，天然符合无监督异常检测设定；
    \item 特征可解释，便于排查异常来自水平偏移、波动增强、快速跳变还是频域振荡；
    \item 使用补全窗口，避免原始缺失值直接破坏均值、标准差和 FFT 计算；
    \item 与原有 \code{detector_score}、\code{disagreement_score} 只通过分数合成耦合，改动面较小。
\end{itemize}

\subsection{当前实现的局限}

\begin{itemize}
    \item 特征分组写死为 \(0{:}6\)、\(6{:}12\)、\(12{:}18\)、\(0{:}18\)，对非 \(18\) 列或不同变量分组的数据不够通用。
    \item 只使用每维标准差，不建模特征之间的相关性，因此不是完整马氏距离。
    \item \(\bm{\mu}_f,\bm{\sigma}_f\) 没有保存到模型文件中。如果训练和测试分开运行，需要额外持久化。
    \item \code{combine_scores} 中没有对 \code{shift score} 的归一化值取绝对值。低于训练均值的偏移会降低最终分数，这是否符合任务目标需要通过实验确认。
    \item \code{shift score} 依赖补全模型。如果补全模型在异常段把异常“抹平”，分布偏移信号也可能被削弱。
\end{itemize}

\section{常见误解与正确理解}

{\small
\begin{longtable}{@{}>{\raggedright\arraybackslash}p{0.32\linewidth} >{\raggedright\arraybackslash}p{0.60\linewidth}@{}}
\toprule
容易误解的说法 & 更准确的理解 \\
\midrule
\endhead
\code{shift score} 是一个训练出来的神经网络分支。 & 不准确。它没有神经网络层，也没有梯度更新。它只是用训练集窗口特征估计均值和标准差。 \\
\code{shift score} 直接比较原始训练集和原始测试集。 & 不准确。它比较的是训练后模型生成的补全窗口，即 \code{train_complete_scaled} 和 \code{test_complete_scaled} 重新切成的窗口。 \\
训练集的 \code{shift score} 越高，说明训练集中有异常。 & 不一定。训练集分数主要用于建立基准分布和阈值。单个训练窗口分数高只能说明它离训练平均统计形态较远。 \\
测试集会重新估计 \code{feature_mean} 和 \code{feature_std}。 & 不会。测试集只使用训练集估计出的 \(\bm{\mu}_f,\bm{\sigma}_f\)，否则会产生测试信息泄漏。 \\
\code{shift score} 越低一定越正常。 & 通常可以这样理解，但最终判断还要看合成分数、EWAF 平滑、阈值和最短异常持续长度。 \\
\bottomrule
\end{longtable}
}

\section{可复现实验与实现建议}

\subsection{建议输出的诊断量}

若后续要进一步研究 \code{shift score}，建议在实验输出中增加以下诊断量：
\begin{itemize}
    \item 每个特征维度的 \(\mu_{f,j}\) 和 \(\sigma_{f,j}\)；
    \item 测试窗口中贡献最大的若干特征维度；
    \item \code{shift_score} 在正常段和异常段的均值、标准差、分位数；
    \item 合成分数前的 \(\widetilde{r}\)、\(\widetilde{g}\)、\(\widetilde{q}\) 三个通道值；
    \item 不同 \code{--score-shift-weight} 下的检测指标变化。
\end{itemize}

这些诊断量可以帮助判断 \code{shift score} 是真的捕捉到了异常分布变化，还是只是被补全模型或尺度因素带偏。

\subsection{如果要把训练和测试拆开}

当前 \code{compute_distribution_shift_scores} 在同一次调用中同时接收训练窗口和测试窗口，因此 \(\bm{\mu}_f,\bm{\sigma}_f\) 是局部变量。若未来需要先训练保存、后单独加载测试，应保存以下内容：
\begin{itemize}
    \item 神经网络权重 \(\bm{\theta}\)；
    \item 标准化器参数 \(\mu_n,\sigma_n\)；
    \item \code{rate_id} 和 \code{stride}；
    \item \code{shift score} 的 \(\bm{\mu}_f,\bm{\sigma}_f\)；
    \item 三个分数通道的训练均值和标准差；
    \item 训练分数阈值 \(\theta\)。
\end{itemize}

否则测试时无法复现训练时的分数尺度。

\subsection{如果要提升通用性}

当前特征构造与 \(18\) 列、三组采样率强绑定。若要支持更多数据集，可以考虑：
\begin{enumerate}
    \item 根据 \code{rate_id} 自动生成变量组，而不是写死列切片；
    \item 把 FFT 频点数写成参数，例如 \code{--shift-fft-bins}；
    \item 对每个特征组输出单独贡献，方便解释；
    \item 试验对 \(\widetilde{q}\) 取绝对值或只保留正偏移：
    \begin{equation}
        \max(\widetilde{q},0),
    \end{equation}
    以避免低于训练均值的分数通道降低总异常分数。
\end{enumerate}

\section{结论}

\code{mra.py} 中的 \code{shift score} 是一个基于训练集补全窗口统计特征的分布偏移分数。它的入口在 \code{compute_distribution_shift_scores}，主流程接入点在 \code{main} 的补全序列重新窗口化之后。它没有梯度训练，也不更新神经网络权重；它从训练集补全窗口中估计特征均值和标准差，再用这些统计量衡量测试补全窗口的偏离程度。

从设计上看，\code{shift score} 与 \code{detector_score} 和 \code{disagreement_score} 互补。前者看补全后的窗口统计形态是否偏离训练分布，后两者分别看检测器重构误差和专家不一致程度。三者经过训练集尺度归一化后合成最终分数，再通过 EWAF、训练阈值和最短异常持续长度得到最终预测。

因此，理解 \code{shift score} 时应抓住三点：第一，它的输入是训练后模型补全出的稠密窗口；第二，它的“训练参数”是训练窗口特征均值和标准差；第三，它在测试集中通过标准化距离和最终分数合成发挥作用，而不是通过神经网络前向层直接输出异常标签。

\end{document}
